\documentclass[conference]{IEEEtran}

\usepackage[utf8]{inputenc}
\usepackage[T1]{fontenc}
\usepackage{amsmath,amssymb}
\usepackage{graphicx}
\usepackage{booktabs}
\usepackage{hyperref}
\usepackage{algorithm}
\usepackage{algpseudocode}
\usepackage{listings}
\usepackage{xcolor}
\usepackage[spanish]{babel}

\hypersetup{
    colorlinks=true,
    linkcolor=blue,
    citecolor=blue,
    urlcolor=blue
}

\title{Destilaci\'on de Modelos de Lenguaje con Generaci\'on Aumentada por Recuperaci\'on para el Dominio Legal Colombiano}

\author{
\IEEEauthorblockN{Equipo de Investigaci\'on}
\IEEEauthorblockA{
Universidad \\
Departamento de Ingenier\'ia de Sistemas \\
Colombia}
}

\begin{document}

\maketitle

\begin{abstract}
Los modelos de lenguaje grandes (LLMs) han demostrado capacidades notables en tareas de comprensi\'on y generaci\'on de texto en dominios especializados. Sin embargo, su alto costo computacional de inferencia limita su despliegue en escenarios con recursos restringidos. Este trabajo presenta un pipeline integral para la destilaci\'on de conocimiento de un modelo Llama-2-7B hacia un modelo TinyLlama-1.1B, aplicado al dominio de documentos legales colombianos. Nuestro enfoque integra Generaci\'on Aumentada por Recuperaci\'on (RAG) durante la fase de entrenamiento del profesor, un sistema de control de calidad para los datos de entrenamiento, y evaluaci\'on autom\'atica mediante LLM-as-Judge. Describimos la arquitectura completa del pipeline ---desde la recolecci\'on y procesamiento de documentos del portal SUIN hasta la evaluaci\'on del modelo destilado--- y discutimos las decisiones t\'ecnicas, resultados preliminares y direcciones futuras de investigaci\'on.
\end{abstract}

\begin{IEEEkeywords}
Destilaci\'on de conocimiento, modelos de lenguaje grandes, RAG, documentos legales, NLP en espa\~nol
\end{IEEEkeywords}

% ============================================================
\section{Introducci\'on}
\label{sec:intro}

El procesamiento de lenguaje natural aplicado al dominio legal ha ganado relevancia debido a la complejidad y volumen de los marcos normativos. En Colombia, el Sistema \'Unico de Informaci\'on Normativa (SUIN) alberga miles de documentos legales cuyo an\'alisis automatizado podr\'ia beneficiar a profesionales del derecho, entidades gubernamentales y ciudadanos.

Los LLMs como GPT-4, Llama-2 y Mistral han demostrado capacidades sobresalientes en comprensi\'on legal \cite{katz2024gpt4}. Sin embargo, desplegar modelos de 7B+ par\'ametros en producci\'on implica costos significativos de infraestructura y latencia elevada, particularmente en contextos donde el procesamiento on-premise es un requisito.

La \textbf{destilaci\'on de conocimiento} \cite{hinton2015distilling} ofrece una soluci\'on: transferir el conocimiento de un modelo grande (\textit{teacher}) a uno m\'as peque\~no (\textit{student}), preservando la mayor parte de las capacidades con una fracci\'on de los par\'ametros. Este trabajo extiende el paradigma de destilaci\'on con dos innovaciones contextuales:

\begin{enumerate}
    \item \textbf{RAG-augmented distillation}: El modelo teacher genera sus respuestas con acceso a un sistema RAG sobre la base legal colombiana, produciendo datos de entrenamiento m\'as precisos y contextualizados.
    \item \textbf{Pipeline integral con control de calidad}: Desde la recolecci\'on de datos hasta la evaluaci\'on, cada etapa incluye mecanismos de validaci\'on y m\'etricas cuantitativas.
\end{enumerate}

% ============================================================
\section{Trabajo Relacionado}
\label{sec:related}

\subsection{Destilaci\'on de Conocimiento en LLMs}

Hinton et al. \cite{hinton2015distilling} propusieron la destilaci\'on de conocimiento usando \textit{soft targets} con un par\'ametro de temperatura. Desde entonces, se han desarrollado m\'ultiples variantes:

\begin{itemize}
    \item \textbf{Basada en logits}: El student aprende a aproximar la distribuci\'on de probabilidades del teacher \cite{hinton2015distilling, kim2016sequence}. Es la t\'ecnica m\'as directa y la que mejor preserva la incertidumbre del teacher.
    \item \textbf{Basada en features}: El student replica representaciones intermedias del teacher \cite{romero2014fitnets, jiao2020tinybert}. Requiere architecturas compatibles.
    \item \textbf{Basada en respuestas}: El student se entrena directamente sobre texto generado por el teacher \cite{taori2023alpaca}. M\'as simple pero pierde informaci\'on distribucional.
\end{itemize}

DistilBERT \cite{sanh2019distilbert} demostr\'o que un modelo 40\% m\'as peque\~no puede retener 97\% del rendimiento del teacher en tareas de comprensi\'on. TinyBERT \cite{jiao2020tinybert} extendi\'o esto con destilaci\'on en m\'ultiples capas.

\subsection{RAG (Retrieval-Augmented Generation)}

Lewis et al. \cite{lewis2020retrieval} introdujeron RAG como un mecanismo para fundamentar las respuestas de LLMs en fuentes externas, reduciendo alucinaciones. En el dominio legal, RAG es particularmente valioso porque las respuestas deben estar ancladas en normas espec\'ificas.

Trabajos recientes han explorado el uso de embeddings multiling\"ues como BGE-M3 \cite{chen2024bge} y t\'ecnicas de reranking de dos etapas \cite{nogueira2019passage} para mejorar la precisi\'on de la recuperaci\'on.

\subsection{LLM como Juez}

Zheng et al. \cite{zheng2023judging} propusieron usar LLMs como jueces autom\'aticos para evaluar calidad de respuestas, demostrando alta correlaci\'on con evaluaci\'on humana. Este enfoque es particularmente \'util cuando la evaluaci\'on manual no escala.

% ============================================================
\section{Metodolog\'ia}
\label{sec:methodology}

\subsection{Arquitectura General del Pipeline}

El pipeline se compone de cinco etapas secuenciales (Figura~\ref{fig:pipeline}):

\begin{enumerate}
    \item Recolecci\'on y preprocesamiento de documentos
    \item Indexaci\'on en base de datos vectorial (RAG)
    \item Generaci\'on de datos de entrenamiento
    \item Destilaci\'on de conocimiento
    \item Evaluaci\'on autom\'atica
\end{enumerate}

\subsection{Recolecci\'on y Preprocesamiento de Datos}
\label{sec:data}

Los documentos legales se obtienen del portal SUIN mediante web scraping con mecanismos de retry y backoff exponencial. Cada documento HTML pasa por un pipeline de limpieza que incluye:

\begin{itemize}
    \item Eliminaci\'on de elementos no textuales (scripts, estilos, navegaci\'on)
    \item Normalizaci\'on de texto legal (40+ reglas regex para manejo de art\'iculos, valores monetarios, listas, y estructura legal)
    \item Extracci\'on de metadatos (tipo, n\'umero, a\~no, entidad emisora, estado)
\end{itemize}

Se implementa un sistema de \textbf{scoring de calidad} con tres dimensiones ponderadas:

\begin{equation}
    Q = w_l \cdot S_{\text{lines}} + w_f \cdot S_{\text{frag}} + w_h \cdot S_{\text{header}}
\end{equation}

donde $S_{\text{lines}}$ es el ratio de l\'ineas cortas (peso 45\%), $S_{\text{frag}}$ la fragmentaci\'on (45\%), y $S_{\text{header}}$ la integridad de encabezados (10\%). Documentos con $Q < 70$ se descartan autom\'aticamente.

\subsection{Sistema RAG}
\label{sec:rag}

Los documentos limpios se indexan en una base de datos vectorial ChromaDB usando el modelo de embeddings BGE-M3 \cite{chen2024bge}, seleccionado por su rendimiento en tareas multiling\"ues y su optimizaci\'on para dominios especializados.

\textbf{Chunking:} Se utiliza \texttt{RecursiveCharacterTextSplitter} con tama\~no de chunk de 800 tokens y overlap de 150 tokens, preservando la coherencia de art\'iculos legales.

\textbf{Recuperaci\'on en dos etapas:}
\begin{enumerate}
    \item B\'usqueda por similitud en ChromaDB ($k=10$)
    \item Reranking con BGE-reranker-v2-m3 ($k_{\text{final}}=2$)
\end{enumerate}

Esta estrategia maximiza recall en la primera etapa y precision en la segunda.

\subsection{Generaci\'on de Datos de Entrenamiento}
\label{sec:datagen}

Se utiliza el modelo teacher (Llama-2-7B-chat) para generar pares de instrucciones-respuestas a partir de los documentos legales. Se definen 9 tipos de tareas de comprensi\'on:

\begin{enumerate}
    \item Idea central
    \item Resumen en 3 niveles
    \item Esencial vs. accesorio
    \item Estructura l\'ogica
    \item Reescritura simplificada
    \item Intenci\'on del autor
    \item Conceptos clave
    \item Reducci\'on extrema
    \item Conexiones internas
\end{enumerate}

Cada documento genera exactamente 10 pares, con validaci\'on de longitud m\'inima (instrucci\'on $\geq 10$ caracteres, respuesta $\geq 30$ caracteres).

\subsection{Destilaci\'on de Conocimiento}
\label{sec:distillation}

Se emplea destilaci\'on basada en logits con una funci\'on de p\'erdida h\'ibrida:

\begin{equation}
    \mathcal{L} = \alpha \cdot \mathcal{L}_{\text{soft}} + (1 - \alpha) \cdot \mathcal{L}_{\text{hard}}
    \label{eq:loss}
\end{equation}

donde:

\begin{equation}
    \mathcal{L}_{\text{soft}} = T^2 \cdot D_{KL}\left(\sigma\left(\frac{z^s}{T}\right) \Big\| \sigma\left(\frac{z^t}{T}\right)\right)
\end{equation}

\begin{equation}
    \mathcal{L}_{\text{hard}} = -\sum_{i} y_i \log \sigma(z^s_i)
\end{equation}

con $T=4.0$ (temperatura), $\alpha=0.7$, $z^s$ y $z^t$ los logits del student y teacher respectivamente, y $\sigma$ la funci\'on softmax.

\textbf{Almacenamiento eficiente de logits:} Dado que la matriz completa de logits del teacher tiene dimensiones $L \times V$ (longitud de secuencia $\times$ tama\~no de vocabulario), almacenarla directamente resulta prohibitivo en t\'erminos de espacio en disco. Para un vocabulario de 32,000 tokens y longitud de secuencia de 2,048, cada muestra generar\'ia $\sim$500\,MB en formato texto. En su lugar, se almacenan \'unicamente los top-$K$ logits ($K=50$) por posici\'on en formato binario (\texttt{.pt}, float16), reduciendo el almacenamiento por muestra de $\sim$500\,MB a $\sim$800\,KB ---una reducci\'on de $\sim$600$\times$---. Durante el entrenamiento, el tensor completo se reconstruye mediante \texttt{scatter\_}, asignando un valor de piso ($-10$) a las posiciones no cubiertas por el top-$K$, lo cual preserva la forma de la distribuci\'on softmax sin afectar significativamente la divergencia KL.

\textbf{Configuraci\'on de entrenamiento:}
\begin{itemize}
    \item Teacher: \texttt{meta-llama/Llama-2-7b-chat-hf} ($\sim$7B par\'ametros)
    \item Student: \texttt{TinyLlama/TinyLlama-1.1B-Chat-v1.0} ($\sim$1.1B par\'ametros)
    \item Optimizer: AdamW, $lr = 2 \times 10^{-5}$
    \item Batch size: 2, \'epocas: 3
    \item Longitud m\'axima de secuencia: 2048 tokens
    \item Top-$K$ logits almacenados: 50
    \item Gradient clipping aplicado
\end{itemize}

\textbf{RAG-augmented distillation:} Opcionalmente, durante la generaci\'on de datos de destilaci\'on, el teacher recibe contexto recuperado del sistema RAG. Esto produce respuestas m\'as fundamentadas en la normativa, que el student aprende a emular.

\begin{algorithm}[t]
\caption{Pipeline de Destilaci\'on con RAG}
\label{alg:pipeline}
\begin{algorithmic}[1]
\State Cargar datos procesados $D = \{(x_i)\}_{i=1}^{N}$
\State Cargar teacher $M_T$ y student $M_S$
\For{cada instrucci\'on $x_i \in D$}
    \If{RAG habilitado}
        \State $c_i \leftarrow \text{Retrieve}(x_i)$ \Comment{Contexto RAG}
        \State $p_i \leftarrow \text{BuildPrompt}(x_i, c_i)$
    \Else
        \State $p_i \leftarrow \text{BuildPrompt}(x_i)$
    \EndIf
    \State $y_i, z_i^t \leftarrow M_T(p_i)$ \Comment{Respuesta + logits}
    \State $\hat{z}_i^t \leftarrow \text{TopK}(z_i^t, K)$ \Comment{Retener top-$K$ logits}
    \State Guardar $(p_i, y_i)$ en JSONL; $\hat{z}_i^t$ en \texttt{.pt}
\EndFor
\For{\'epoca $= 1$ to $E$}
    \For{cada batch $B$ del dataset}
        \State $z^t \leftarrow \text{Reconstruct}(\hat{z}^t, V)$ \Comment{Expandir a vocabulario completo}
        \State $z^s \leftarrow M_S(B)$ \Comment{Logits del student}
        \State $\mathcal{L} \leftarrow \alpha \cdot \mathcal{L}_{\text{soft}} + (1-\alpha) \cdot \mathcal{L}_{\text{hard}}$
        \State Actualizar $M_S$ con $\nabla \mathcal{L}$
    \EndFor
\EndFor
\end{algorithmic}
\end{algorithm}

\subsection{Evaluaci\'on con LLM-as-Judge}
\label{sec:evaluation}

El modelo destilado se eval\'ua mediante un framework de LLM-as-Judge que puntua las respuestas en cuatro dimensiones:

\begin{table}[h]
\centering
\caption{Dimensiones de evaluaci\'on y pesos}
\label{tab:eval_dimensions}
\begin{tabular}{lcc}
\toprule
\textbf{Dimensi\'on} & \textbf{Peso} & \textbf{Escala} \\
\midrule
Accuracy (precisi\'on factual) & 40\% & 1--5 \\
Relevance (pertinencia) & 30\% & 1--5 \\
Completeness (cobertura) & 20\% & 1--5 \\
Clarity (claridad) & 10\% & 1--5 \\
\bottomrule
\end{tabular}
\end{table}

La puntuaci\'on ponderada final se calcula como:

\begin{equation}
    S = 0.4 \cdot S_a + 0.3 \cdot S_r + 0.2 \cdot S_c + 0.1 \cdot S_{cl}
\end{equation}

Adicionalmente, se eval\'ua la calidad del sistema RAG usando m\'etricas est\'andar de recuperaci\'on:
\begin{itemize}
    \item Contextual Relevancy
    \item Contextual Recall
    \item Contextual Precision
\end{itemize}

% ============================================================
\section{Resultados Preliminares}
\label{sec:results}

\subsection{Recolecci\'on de Datos}

El pipeline de recolecci\'on proces\'o documentos del portal SUIN, clasificando cada archivo seg\'un su puntaje de calidad. Los documentos con puntaje $\geq 70$ se consideran usables para entrenamiento.

\subsection{Datos de Entrenamiento}

Se generaron pares de instruccion-respuesta cubriendo los 9 tipos de tareas de comprensi\'on definidos, proporcionando diversidad en el entrenamiento del modelo student.

\textit{Nota: Los resultados cuantitativos completos de destilaci\'on y evaluaci\'on se reportar\'an en la versi\'on final del art\'iculo una vez completados los experimentos con los baselines de comparaci\'on.}

% ============================================================
\section{Discusi\'on}
\label{sec:discussion}

\subsection{Ventajas del Enfoque}

\begin{itemize}
    \item \textbf{Pipeline integral}: A diferencia de trabajos que se enfocan solo en la destilaci\'on, nuestro enfoque cubre todo el ciclo de vida: desde la recolecci\'on de datos hasta la evaluaci\'on.
    \item \textbf{Control de calidad en datos}: El sistema de scoring autom\'atico reduce el ruido en los datos de entrenamiento.
    \item \textbf{RAG durante destilaci\'on}: El teacher produce respuestas m\'as fundamentadas, lo que potencialmente mejora la calidad del student.
    \item \textbf{Dominio espec\'ifico}: El enfoque en documentos legales colombianos demuestra aplicabilidad a idiomas y dominios poco representados en la literatura.
\end{itemize}

\subsection{Limitaciones}

\begin{itemize}
    \item \textbf{Tama\~no del student}: TinyLlama-1.1B puede ser insuficiente para capturar la complejidad del lenguaje legal. Modelos intermedios (2-3B par\'ametros) podr\'ian ofrecer mejor balance.
    \item \textbf{Teacher no \'optimo}: Llama-2-7B no es el modelo m\'as capaz disponible. Modelos m\'as recientes (Llama-3, Mistral) como teachers podr\'ian mejorar la calidad de la destilaci\'on.
    \item \textbf{Evaluaci\'on humana pendiente}: La evaluaci\'on por LLM-as-Judge debe validarse contra juicio de expertos legales.
    \item \textbf{Effective batch size}: El batch size de 2 sin gradient accumulation puede afectar la convergencia.
\end{itemize}

\subsection{Direcciones Futuras}

\begin{enumerate}
    \item Evaluar modelos student de mayor capacidad (Phi-2, Mistral-7B con cuantizaci\'on)
    \item Implementar destilaci\'on progresiva (7B $\rightarrow$ 3B $\rightarrow$ 1B)
    \item Explorar destilaci\'on basada en features adem\'as de logits
    \item Comparar con fine-tuning directo del student (sin destilaci\'on)
    \item Validar con evaluadores humanos expertos en derecho colombiano
    \item Explorar cuantizaci\'on (GPTQ, AWQ) como alternativa complementaria
\end{enumerate}

% ============================================================
\section{Conclusi\'on}
\label{sec:conclusion}

Este trabajo presenta un pipeline completo para la destilaci\'on de LLMs en el dominio legal colombiano, integrando RAG, control de calidad de datos, y evaluaci\'on autom\'atica. La arquitectura modular permite la iteraci\'on independiente sobre cada componente, facilitando la experimentaci\'on y la mejora continua. Los resultados preliminares demuestran la viabilidad del enfoque, y los experimentos futuros con baselines de comparaci\'on permitir\'an cuantificar el aporte espec\'ifico de cada componente del pipeline.

% ============================================================
\section*{Agradecimientos}

Los autores agradecen el acceso a recursos computacionales de HPC (Apolo) y al portal SUIN por la disponibilidad p\'ublica de documentos legales colombianos.

% ============================================================
\bibliographystyle{IEEEtran}

\begin{thebibliography}{10}

\bibitem{hinton2015distilling}
G.~Hinton, O.~Vinyals, and J.~Dean, ``Distilling the knowledge in a neural network,'' \textit{arXiv preprint arXiv:1503.02531}, 2015.

\bibitem{sanh2019distilbert}
V.~Sanh, L.~Debut, J.~Chaumond, and T.~Wolf, ``DistilBERT, a distilled version of BERT: smaller, faster, cheaper and lighter,'' \textit{arXiv preprint arXiv:1910.01108}, 2019.

\bibitem{jiao2020tinybert}
X.~Jiao \textit{et al.}, ``TinyBERT: Distilling BERT for natural language understanding,'' in \textit{Findings of EMNLP}, 2020, pp.~4163--4174.

\bibitem{lewis2020retrieval}
P.~Lewis \textit{et al.}, ``Retrieval-augmented generation for knowledge-intensive NLP tasks,'' in \textit{NeurIPS}, vol.~33, 2020, pp.~9459--9474.

\bibitem{chen2024bge}
J.~Chen \textit{et al.}, ``BGE M3-Embedding: Multi-lingual, multi-functionality, multi-granularity text embeddings through self-knowledge distillation,'' \textit{arXiv preprint arXiv:2402.03216}, 2024.

\bibitem{zheng2023judging}
L.~Zheng \textit{et al.}, ``Judging LLM-as-a-Judge with MT-Bench and Chatbot Arena,'' in \textit{NeurIPS}, vol.~36, 2023.

\bibitem{katz2024gpt4}
D.~M.~Katz \textit{et al.}, ``GPT-4 passes the bar exam,'' \textit{Philosophical Transactions of the Royal Society A}, vol.~382, no.~2270, 2024.

\bibitem{kim2016sequence}
Y.~Kim and A.~M.~Rush, ``Sequence-level knowledge distillation,'' in \textit{EMNLP}, 2016, pp.~1317--1327.

\bibitem{romero2014fitnets}
A.~Romero \textit{et al.}, ``FitNets: Hints for thin deep nets,'' in \textit{ICLR}, 2015.

\bibitem{taori2023alpaca}
R.~Taori \textit{et al.}, ``Stanford Alpaca: An instruction-following LLaMA model,'' 2023.

\bibitem{nogueira2019passage}
R.~Nogueira and K.~Cho, ``Passage re-ranking with BERT,'' \textit{arXiv preprint arXiv:1901.04085}, 2019.

\end{thebibliography}

\end{document}
